\documentclass[11pt,a4paper]{report}

% --- Pacchetti Fondamentali ---
\usepackage[utf8]{inputenc}
\usepackage[T1]{fontenc}
\usepackage{lmodern}
\usepackage[italian]{babel}
\usepackage{etoolbox}
\usepackage[margin=1in]{geometry}
\usepackage{hyperref}
\usepackage{xcolor}
\usepackage{listings}
\usepackage{enumitem}
\usepackage{titlesec}
\usepackage{fancyhdr}
\usepackage{hyperref}
\usepackage{cleveref}

\crefformat{section}{\S#2#1#3}
\crefformat{subsection}{\S#2#1#3}
\crefformat{subsubsection}{\S#2#1#3}
\crefformat{figure}{#2Figura~#1#3}
\crefformat{footnote}{#2\footnotemark[#1]#3}
\crefformat{table}{#2Tabella~#1#3}

% --- Configurazione Header e Footer ---
\pagestyle{fancy}
\fancyhf{}
\fancyhead[L]{\small Programmazione e Tecnologie Web}
\fancyhead[R]{\small Giuseppe Capasso}
\fancyfoot[C]{\thepage}
\renewcommand{\headrulewidth}{0.4pt}
\usepackage{tikz}
\usetikzlibrary{arrows.meta, positioning, shapes.geometric, decorations.pathreplacing}

% --- Configurazione Colori e Link ---
\definecolor{primary}{RGB}{38, 66, 139} % Blu scuro istituzionale
\definecolor{codebg}{RGB}{245, 245, 245}

\hypersetup{
    colorlinks=true,
    linkcolor=primary,
    urlcolor=primary,
    citecolor=primary,
}

% --- Configurazione Formattazione Codice ---
\lstset{
    backgroundcolor=\color{codebg},
    basicstyle=\ttfamily\small,
    breaklines=true,
    frame=single,
    rulecolor=\color{lightgray},
    captionpos=b,
    xleftmargin=1em,
    xrightmargin=1em
}

% --- Formattazione Titoli ---
\titleformat{\chapter}[display]
  {\normalfont\bfseries\color{primary}}{\huge Capitolo \thechapter}{1ex}{\Huge}
\titlespacing*{\chapter}{0pt}{-20pt}{40pt}

\begin{document}

% --- Frontespizio ---
\begin{titlepage}
    \centering
    \vspace*{4cm}
    {\Huge \textbf{\color{primary} Programmazione e Tecnologie Web}} \\[1.5cm]
    {\Large \textbf{Dispense del Corso - Lezione 2}} \\[0.5cm]
    {\large Docente: Giuseppe Capasso} \\[1.5cm]
    {\normalsize \textbf{Cloud Native Development \& Cybersecurity Specialist}} \\
    {\normalsize Istituto Tecnico Superiore (ITS)} \\
    \vfill
    {\normalsize Anno Accademico 2025/2026}
\end{titlepage}

\newpage

% ==========================================
% INTRODUZIONE
% ==========================================
\chapter*{Il Web Statico vs Dinamico e la Logica Client-Side}
\addcontentsline{toc}{chapter}{Il Web Statico vs Dinamico e la Logica Client-Side}

\section*{Obiettivi della Lezione}
In questa lezione esploreremo l'evoluzione delle architetture web, passando dalla semplicita' dei siti statici alla complessita' delle moderne applicazioni web dinamiche. Gli obiettivi principali sono:
\begin{itemize}
    \item \textbf{Distinguere tra Web Statico e Dinamico}: Comprendere quando utilizzare l'uno o l'altro approccio in base alle esigenze del progetto.
    \item \textbf{Paradigmi di Rendering}: Analizzare le differenze tra Server-Side Rendering (SSR) e Client-Side Rendering (CSR).
    \item \textbf{Generazione di Siti Statici (SSG)}: Imparare a utilizzare \textbf{Hugo} per creare siti veloci e sicuri.
    \item \textbf{Persistenza dei Dati}: Gestire lo stato dell'applicazione direttamente nel browser utilizzando il \textbf{localStorage}.
    \item \textbf{Esercitazione Pratica}: Implementare una Todo App completa seguendo il paradigma CRUD.
\end{itemize}

% ==========================================
% CAPITOLO 1
% ==========================================
\chapter{Web Statico vs Web Dinamico}

\section{Il Web Statico: Il Ritorno alle Origini}
Un sito web si definisce \textbf{statico} quando il contenuto servito al client e' esattamente lo stesso memorizzato sul file system del server. Non vi e' alcuna elaborazione in tempo reale: il server web (come Apache o Nginx) agisce semplicemente come un distributore di file.

Quando un utente richiede una pagina (es. \texttt{contatti.html}), il server cerca il file corrispondente e lo invia nel corpo della risposta HTTP.

\begin{figure}[ht]
\centering
\begin{tikzpicture}[node distance=4cm, every node/.style={rectangle, draw, rounded corners, minimum width=2.5cm, minimum height=1.2cm, align=center, line width=1pt}]
    \node (client) [fill=primary!10] {\textbf{Client} \\ (Browser)};
    \node (server) [right=of client, fill=primary!10] {\textbf{Server Web} \\ (Nginx/Apache)};
    \node (disk) [below=1.5cm of server, shape=cylinder, aspect=0.5, rotate=90, minimum width=1.5cm, minimum height=2cm, fill=gray!20] {\rotatebox{-90}{File System}};

    \draw [->, >=stealth, line width=1.2pt, primary] ([yshift=0.3cm]client.east) -- node[above, black, scale=0.8] {GET /index.html} ([yshift=0.3cm]server.west);
    \draw [->, >=stealth, line width=1pt, dashed] (server.south) -- node[right, black, scale=0.8] {Leggi file} (disk.north);
    \draw [<-, >=stealth, line width=1.2pt, primary] ([yshift=-0.3cm]client.east) -- node[below, black, scale=0.8] {200 OK (index.html)} ([yshift=-0.3cm]server.west);
\end{tikzpicture}
\caption{Architettura di un sito web statico.}
\label{fig:static_web}
\end{figure}

\subsection{Vantaggi del Web Statico}
Nonostante possa sembrare una tecnologia superata, il web statico sta vivendo una rinascita grazie a:
\begin{itemize}
    \item \textbf{Performance Estreme}: Poiche' non c'e' elaborazione lato server (nessuna query al database, nessuna esecuzione di codice PHP/Python), la risposta e' quasi istantanea.
    \item \textbf{Sicurezza}: Senza un database o un backend attivo, la superficie d'attacco e' minima. Non e' possibile effettuare SQL Injection o altri attacchi comuni ai CMS dinamici come WordPress.
    \item \textbf{Scalabilita'}: Servire file statici richiede pochissime risorse. Una CDN (Content Delivery Network) puo' replicare il sito in tutto il mondo con costi irrisori.
\end{itemize}

\subsection{Svantaggi e Limiti}
Il limite principale e' la \textbf{gestione dei contenuti}. In un sito statico "puro", se vuoi cambiare il footer su 100 pagine, devi modificare manualmente 100 file HTML. Inoltre, non e' possibile avere contenuti personalizzati per l'utente (es. "Benvenuto Giuseppe") o gestire login e carrelli acquisti senza l'ausilio di servizi esterni o JavaScript.

\section{Il Web Dinamico: Personalizzazione e Interazione}
Un sito web e' \textbf{dinamico} quando il contenuto della pagina viene generato "al volo" dal server in risposta a una specifica richiesta. In questo modello, il server non si limita a leggere un file, ma esegue del codice (in PHP, Python, JavaScript, etc.) che spesso interagisce con un \textbf{Database}.

\begin{figure}[ht]
\centering
\begin{tikzpicture}[node distance=3cm, every node/.style={rectangle, draw, rounded corners, minimum width=2.2cm, minimum height=1cm, align=center, line width=1pt}]
    \node (client) [fill=primary!10] {\textbf{Client}};
    \node (server) [right=of client, fill=primary!20] {\textbf{App Server} \\ (Backend)};
    \node (db) [right=of server, shape=cylinder, aspect=0.5, rotate=90, minimum width=1.2cm, minimum height=1.5cm, fill=primary!5] {\rotatebox{-90}{Database}};

    \draw [->, >=stealth, line width=1.2pt, primary] ([yshift=0.2cm]client.east) -- node[above, black, scale=0.7] {Request} ([yshift=0.2cm]server.west);
    \draw [<->, >=stealth, line width=1pt] (server.east) -- node[above, black, scale=0.7] {Query} (db.north);
    \draw [<-, >=stealth, line width=1.2pt, primary] ([yshift=-0.2cm]client.east) -- node[below, black, scale=0.7] {HTML Generato} ([yshift=-0.2cm]server.west);
\end{tikzpicture}
\caption{Architettura di un sito web dinamico tradizionale.}
\label{fig:dynamic_web}
\end{figure}

Il vantaggio principale e' la possibilita' di gestire enormi quantita' di dati e offrire un'esperienza utente su misura. Tuttavia, questo comporta una maggiore complessita' infrastrutturale e potenziali colli di bottiglia nelle prestazioni dovuti alle query al database.

\section{Esempio Pratico: Hosting Gratuito con Cloudflare Pages}
Per mettere in pratica i vantaggi del web statico, oggi esistono servizi che permettono di ospitare siti web in modo professionale e totalmente gratuito. Uno dei piu' famosi e' \textbf{Cloudflare Pages}.

\subsection{Perche' usare una CDN (Content Delivery Network)?}
In un modello tradizionale, il tuo sito vive su un unico server in una specifica posizione geografica (es. Milano). Se un utente da Tokyo cerca di accedere, i dati devono attraversare l'intero globo, introducendo latenza.

Una \textbf{CDN} come Cloudflare risolve questo problema distribuendo copie del tuo sito statico su centinaia di server (chiamati \textit{Edge Nodes}) in tutto il mondo. L'utente verra' servito dal server fisicamente piu' vicino a lui.

\begin{figure}[ht]
\centering
\begin{tikzpicture}[node distance=1.5cm, every node/.style={align=center, font=\small}]
    \node (origin) [rectangle, draw, fill=primary!20, minimum width=2.5cm, minimum height=1cm] {\textbf{Server Origine}};
    
    \node (edge1) [circle, draw, fill=primary!5, below left=1cm and 2cm of origin] {Edge\\Londra};
    \node (edge2) [circle, draw, fill=primary!5, below=1.5cm of origin] {Edge\\Milano};
    \node (edge3) [circle, draw, fill=primary!5, below right=1cm and 2cm of origin] {Edge\\Tokyo};

    \draw [->, >=stealth, dashed] (origin) -- (edge1);
    \draw [->, >=stealth, dashed] (origin) -- (edge2);
    \draw [->, >=stealth, dashed] (origin) -- (edge3);

    \node (user) [rectangle, draw, rounded corners, fill=green!10, below=1cm of edge3] {Utente Japan};
    \draw [<->, >=stealth, primary, line width=1pt] (edge3) -- (user);
\end{tikzpicture}
\caption{Funzionamento di una CDN: i contenuti statici sono vicini all'utente.}
\label{fig:cdn_working}
\end{figure}

\subsection{Workflow con Cloudflare Pages}
Il "trucco" per ospitare un sito statico gratuitamente su Cloudflare Pages consiste nell'integrazione con Git:
\begin{enumerate}
    \item \textbf{Codice su GitHub}: Carichi i tuoi file HTML/CSS/JS su un repository GitHub.
    \item \textbf{Collegamento}: Connetti il tuo account Cloudflare al repository.
    \item \textbf{Deployment Automatico}: Ogni volta che fai un \texttt{git push}, Cloudflare rileva la modifica, preleva i file e li distribuisce sulla sua rete globale in pochi secondi.
\end{enumerate}

Questo approccio non solo e' gratuito per la maggior parte dei progetti personali, ma garantisce anche certificati SSL (HTTPS) automatici e una velocita' di caricamento imbattibile.

\subsection{Mini-Tutorial: Pubblicare il tuo primo sito statico}
Segui questi passaggi per rendere il tuo progetto accessibile a chiunque su Internet:

\begin{enumerate}
    \item \textbf{Prepara il codice}: Assicurati che il tuo progetto abbia un file chiamato \texttt{index.html} nella cartella principale.
    \item \textbf{Carica su GitHub}:
    \begin{itemize}
        \item Vai su \url{https://github.com} e crea un nuovo repository (es. "mio-sito-statico").
        \item Carica i tuoi file usando i comandi Git visti nella Lezione 1: \\
        \texttt{git add .}, \texttt{git commit -m "Primo commit"}, \texttt{git push}.
    \end{itemize}
    \item \textbf{Configura Cloudflare}:
    \begin{itemize}
        \item Vai sulla dashboard di Cloudflare (\url{https://dash.cloudflare.com/sign-up}) e crea un account gratuito.
        \item Nel menu laterale, clicca su \textbf{Workers \& Pages} e poi sul pulsante \textbf{Create application}.
        \item Seleziona la scheda \textbf{Pages} e clicca su \textbf{Connect to Git}.
    \end{itemize}
    \item \textbf{Deployment}:
    \begin{itemize}
        \item Seleziona il tuo repository GitHub appena creato.
        \item Nelle impostazioni di build, se il tuo sito e' in HTML semplice, puoi lasciare tutto invariato (Production branch: \texttt{main}, Build command: vuoto).
        \item Clicca su \textbf{Save and Deploy}.
    \end{itemize}
\end{enumerate}

Dopo pochi secondi, Cloudflare ti fornira' un indirizzo del tipo \texttt{mio-sito.pages.dev}. Da questo momento, ogni volta che modificherai il codice e farai un \texttt{push} su GitHub, il tuo sito si aggiornera' automaticamente in tutto il mondo!

% ==========================================
% CAPITOLO 2
% ==========================================
\chapter{Paradigmi di Rendering: SSR vs CSR}

\section{Introduzione ai Modelli di Rendering}
La differenza fondamentale tra i vari tipi di siti web risiede in \textbf{dove} viene costruito l'HTML che l'utente vede. Esistono due approcci principali:

\begin{itemize}
    \item \textbf{Server-Side Rendering (SSR)}: La pagina viene "assemblata" sul server. Il browser riceve un file HTML gia' completo di dati.
    \item \textbf{Client-Side Rendering (CSR)}: Il server invia un file HTML quasi vuoto e un grosso file JavaScript. Il browser esegue il JavaScript, che scarica i dati e costruisce l'HTML direttamente nel browser.
\end{itemize}

\begin{figure}[ht]
\centering
\begin{tikzpicture}[node distance=3.5cm, every node/.style={rectangle, draw, rounded corners, minimum width=2.5cm, minimum height=1cm, align=center, line width=1pt}]
    % SSR
    \node (s_server) [fill=primary!20] {\textbf{Server (SSR)}};
    \node (s_client) [right=of s_server, fill=primary!10] {\textbf{Client}};
    \draw [->, >=stealth, primary, line width=1.2pt] (s_server) -- node[above, black, scale=0.7] {HTML Completo\\(Dati inclusi)} (s_client);
    \node at (s_server.north) [above=5pt, font=\bfseries] {Modello Tradizionale};

    % CSR
    \node (c_server) [below=1.5cm of s_server, fill=primary!20] {\textbf{Server (CSR)}};
    \node (c_client) [right=of c_server, fill=primary!10] {\textbf{Client}};
    \draw [->, >=stealth, primary, line width=1.2pt] ([yshift=0.2cm]c_server.east) -- node[above, black, scale=0.7] {HTML Vuoto + JS} ([yshift=0.2cm]c_client.west);
    \draw [->, >=stealth, gray, dashed] ([yshift=-0.2cm]c_client.west) -- node[below, black, scale=0.7] {Richiesta Dati (JSON)} ([yshift=-0.2cm]c_server.east);
    \node at (c_server.north) [above=5pt, font=\bfseries] {Modello Moderno (SPA)};
\end{tikzpicture}
\caption{Confronto tra Server-Side Rendering e Client-Side Rendering.}
\label{fig:ssr_vs_csr}
\end{figure}

\section{SSR e il Motore di Templating}
Nelle applicazioni SSR, utilizziamo un \textbf{Motore di Templating}. Si tratta di un sistema che permette di mescolare codice HTML statico con dei "segnaposto" (placeholder) che verranno sostituiti dai dati reali prima che la pagina venga inviata.

\subsection{Esempio con Django (Python)}
Django utilizza un linguaggio di template molto potente. Immaginiamo di voler mostrare una lista di task salvati in un database Python:

\begin{lstlisting}[language=Python, caption=Esempio di Template Django]
<!-- template.html -->
<h1>Le tue Cose da Fare</h1>
<ul>
    {% for task in lista_todo %}
        <li>
            {{ task.titolo }} 
            {% if task.completato %}
                <span style="color: green;">(Fatto!)</span>
            {% endif %}
        </li>
    {% empty %}
        <li>Non hai ancora impegni.</li>
    {% endfor %}
</ul>
\end{lstlisting}

\textbf{Cosa succede?}: Il server Python legge questo file, esegue il ciclo \texttt{for}, inietta i nomi dei task e invia al browser un \texttt{<ul>} standard. Il browser non sapra' mai che quel file e' stato generato da Python!

\section{Static Site Generation (SSG): Hugo}
Un caso particolare di SSR e' lo \textbf{Static Site Generator}. Invece di generare la pagina ogni volta che un utente la richiede, la generiamo \textbf{una sola volta} durante la fase di sviluppo (build).

\textbf{Hugo} e' uno degli strumenti piu' popolari per questo scopo. Scrivi i contenuti in \textbf{Markdown} (un formato di testo semplice) e Hugo li trasforma in un intero sito web statico in pochi millisecondi.

\begin{lstlisting}[caption=Esempio di contenuto Hugo (Markdown)]
---
title: "Il mio primo post"
date: 2026-05-14
---
# Benvenuti!
Questo e' un esempio di contenuto che Hugo trasformera' in un file HTML statico utilizzando un template.
\end{lstlisting}

% ==========================================
% CAPITOLO 3
% ==========================================
\chapter{Laboratorio: Todo App Client-Side}

In questo laboratorio costruiremo una \textbf{Todo Application} funzionale che gira interamente nel browser. Useremo questo esercizio per comprendere come gestire i dati (\textbf{CRUD}) e come farli sopravvivere al ricaricamento della pagina.

\section{Fase 1: Struttura e Stile (HTML/CSS)}
Iniziamo creando un'interfaccia pulita. Useremo un campo di testo per inserire i task e una lista non ordinata per visualizzarli.

\begin{lstlisting}[language=HTML, caption=Struttura base Todo App]
<div class="container">
    <h1>I miei Task</h1>
    <div class="input-group">
        <input type="text" id="todoInput" placeholder="Cosa devi fare?">
        <button id="addBtn">Aggiungi</button>
    </div>
    <ul id="todoList">
        <!-- I task appariranno qui -->
    </ul>
</div>
\end{lstlisting}

\section{Fase 2: Introduzione a jQuery}
Sebbene JavaScript moderno (Vanilla JS) sia molto potente, per anni \textbf{jQuery} e' stato lo standard per manipolare il DOM in modo semplice e cross-browser. Lo useremo oggi per la sua sintassi concisa.

Per includerlo, aggiungiamo questo tag prima della fine del \texttt{<body>}:
\begin{lstlisting}[language=HTML]
<script src="https://code.jquery.com/jquery-3.7.1.min.js"></script>
\end{lstlisting}

La sintassi di base e' \texttt{\$(selettore).azione()}. Ad esempio:
\begin{lstlisting}[language=Java]
// Cerca il bottone e aggiunge un evento click
$('#addBtn').click(function() {
    alert("Hai cliccato!");
});
\end{lstlisting}

\section{Fase 3: Implementazione del Paradigma CRUD}
Il CRUD rappresenta le quattro operazioni fondamentali:
\begin{enumerate}
    \item \textbf{Create}: Quando premiamo "Aggiungi", creiamo un nuovo oggetto task e lo aggiungiamo a un array.
    \item \textbf{Read}: Leggiamo l'array e generiamo i tag \texttt{<li>} nel DOM.
    \item \textbf{Update}: Cliccando su un task, cambiamo il suo stato (es. "completato: true").
    \item \textbf{Delete}: Un tasto "X" per rimuovere il task dall'array e dalla vista.
\end{enumerate}

\section{Fase 4: Persistenza con localStorage}
Il problema di JavaScript e' che se ricarichi la pagina, tutte le variabili vengono resettate. Per salvare i dati, usiamo il \textbf{localStorage}.

Il \texttt{localStorage} puo' salvare solo \textbf{stringhe}. Per salvare oggetti complessi (come la nostra lista di task), dobbiamo convertirli in formato \textbf{JSON}.

\begin{lstlisting}[language=Java, caption=Salvataggio e Caricamento dati]
// Trasformiamo l'array in stringa e salviamo
const datiString = JSON.stringify(miaListaTasks);
localStorage.setItem('tasks', datiString);

// Recuperiamo la stringa e la ritrasformiamo in oggetto JS
const recuperati = localStorage.getItem('tasks');
const listaOriginale = JSON.parse(recuperati);
\end{lstlisting}

\paragraph{Esercizio Pratico}
Implementa la funzione di salvataggio ogni volta che un task viene aggiunto o rimosso. All'avvio della pagina, controlla se ci sono dati nel \texttt{localStorage} e, se presenti, popolala automaticamente.

\end{document}
